\documentclass[12pt]{article}
\usepackage[utf8]{inputenc}
\usepackage[T1]{fontenc}
\usepackage{hyperref}
\title{Satawalese Language Resources for K-12 Curriculum Development}
\author{Compiled by Claude Sonnet 5.5 (claude-sonnet-5-5) for Andr\'as Kornai}
\date{}
\begin{document}
\maketitle

\section{Introduction}
Satawalese (ISO 639-3 \texttt{stw}, Glottocode sata1237) is a Chuukic (Austronesian) language of Satawal Island, Yap State, Micronesia. Estimates range from about 460 to 700 speakers. Web search found few resources.

\section{Primary Educational Resources}
No specific primers or textbooks were found. A Wikipedia mirror says alphabet books, translations and dictionaries exist, but it names none of them.

\section{Linguistic and Reference Materials}
\begin{itemize}
\item Kevin Roddy, \emph{A Sketch Grammar of Satawalese, the Language of Satawal Island, Yap State, Micronesia} (University of Hawaii thesis, 2007): \url{https://habeleinstitute.org/wiki/Zotero_reference:LS24IPF6}
\item Wikipedia, Satawalese language: \url{https://en.wikipedia.org/wiki/Satawalese_language}
\item Numeral data from the Max Planck Institute for Evolutionary Anthropology: \url{https://lingweb.eva.mpg.de/channumerals/Satawalese.htm}
\item Wiktionary entry for a Satawalese verb, citing Roddy: \url{https://en.wiktionary.com/wiki/nikiti}
\item Verbix Satawalese verb list: \url{https://docs.verbix.com/Languages/Satawalese}
\end{itemize}
No standalone dictionary was confirmed.

\section{Religious and Community Resources}
None found.

\section{Digital and Audio Resources}
\begin{itemize}
\item Endangered Languages Project entry: \url{https://endangeredlanguages.com/lang/5426}
\end{itemize}

\section{Curriculum Development Context}
No evidence of Satawalese-medium schooling was found. The sketch grammar is the best linguistic base for orthography and materials work. A curriculum would need new materials developed with the community.

\end{document}